% packages to import
\usepackage[utf8]{inputenc}
\usepackage[english,ngerman]{babel}
\usepackage{array,tikz,listings,color,graphicx,csquotes,libertine,datetime,enumitem,amssymb}
%hypertexnames=false is needed for complex document structure (more dumb, but more robust)
\usepackage[hypertexnames=false]{hyperref}
\usetikzlibrary{shapes.misc} 
\usepackage{pgfplots} \pgfplotsset{compat=1.14}
\usepackage{wrapfig,tikzscale}

\usepackage{colortbl}


%change the way of numbering stuff:
\numberwithin{equation}{section}
\numberwithin{table}{section}
\numberwithin{figure}{section}
% restart section numbering in every part (english/german)
\makeatletter
\@addtoreset{section}{part} 
\@addtoreset{equation}{part}
\@addtoreset{figure}{part}
\@addtoreset{table}{part}
\@addtoreset{footnote}{section}
\makeatother
% every language gets its own invisiblepart:
\newcommand\invisiblepart[1]{%
  \refstepcounter{part}%
  \addcontentsline{toc}{part}{#1}%
}
% with this one can have multiple \maketitle commands.
\makeatletter
\let\@cleartopmattertags\relax
\makeatother

%two commands for german and english parts of document
\newif\ifen% entrue: english; enfalse: german
\newcommand{\en}[1]{\ifen#1\fi} % print english text only if \entrue
\newcommand{\de}[1]{\ifen\else#1\fi} % print german text only if \enfalse

% layout of Definitions, Theorems, etc.
\newtheorem{defi}{\en{Definition}\de{Definition}}[section]
\newtheorem{lem}[defi]{\en{Lemma}\de{Lemma}}
\newtheorem{thm}[defi]{\en{Proposition}\de{Satz}}
\newtheorem{verm}[defi]{\en{Conjecture}\de{Vermutung}}
\newtheorem{bem}[defi]{\en{Remark}\de{Bemerkung}}
\usepackage{mdframed}
 	\definecolor{lightergray}{rgb}{0.9, 0.9, 0.9}
\newmdtheoremenv[innertopmargin=0pt, backgroundcolor=lightergray, linecolor=gray, linewidth=1pt, topline=false, bottomline=false]{theo}[defi]{\en{Theorem}\de{Theorem}}

%pagestyle, headings:
\usepackage{fancyhdr}
\setlength{\headheight}{12.37082pt}
\pagestyle{fancy}
\fancyhf{}
\fancyhead[OR,EL]{--~\thepage~--} % Odd Right, Even Left
\fancyhead[ER]{\emph{L.~Milla: \en{A Detailed Proof of the Chudnovsky Formula}\de{Ein ausführlicher Beweis der Chudnovsky-Formel}}}
\fancyhead[OL]{\emph{\thesection.~\leftmark}}
\renewcommand{\headrulewidth}{0.4pt}

%change layout of section headings
\makeatletter
  \def\section{\@startsection{section}{1}%
    \z@{.7\linespacing\@plus\linespacing}{.6\linespacing}%
    {\Large\normalfont\scshape\bfseries\centering}}
\makeatother

% individual commands:
\newcommand{\ko}{\en{.}\de{{,}}} % decimal separator english/german
\newcommand{\im}{\operatorname{Im}} % imaginary part
\newcommand{\dsf}[2]{#1/#2} %\frac{}{} written in one line
\DeclareMathOperator{\DIV}{DIV} 
% Cases: Odd/Even
\newcommand{\cas}[2]{\begin{cases}#1 &\text{ \en{if}\de{wenn} } m \text{ \en{is even}\de{gerade ist}}\\%
                     #2 &\text{ \en{if}\de{wenn} } m \text{ \en{is odd}\de{ungerade ist}}\end{cases}}
%used in Table of Contents:
\newcommand{\kommentInh}[2]{{\noindent\textbf{\ref{\en{EN}#1}.~~~\nameref{\en{EN}#1}~~\dotfill~~\pageref{\en{EN}#1}}\par\vspace*{1.5pt}
                            \noindent\narrower\narrower{\itshape #2 }\par\vspace*{7pt}}}
%the following command deletes the extra spacing after delta symbol:
\newcommand{\dl}[1]{\delta\mathopen{}\left(#1\right)\mathclose{}}
%\left( and \right) without spacing
\def\lk{\mathopen{}\left(}
\def\rk{\right)\mathclose{}}


